\section{Определение основной и расширенной матрицы системы. Формулировка теоремы Кронекера--Капелли}

\begin{definition}[Основная и расширенная матрицы системы]
\textbf{Основная матрица системы} --- матрица коэффициентов при неизвестных:
\[
A = \begin{pmatrix}
a_{11} & \cdots & a_{1n} \\
\vdots & \ddots & \vdots \\
a_{m1} & \cdots & a_{mn}
\end{pmatrix}_{m \times n}.
\]
\textbf{Расширенная матрица системы} --- матрица, дополненная столбцом свободных членов:
\[
\tilde{A} = \begin{pmatrix}
a_{11} & \cdots & a_{1n} & b_1 \\
\vdots & \ddots & \vdots & \vdots \\
a_{m1} & \cdots & a_{mn} & b_m
\end{pmatrix}_{m \times (n+1)}.
\]
\end{definition}

\begin{theorem}[Кронекера -- Капелли]
СЛАУ совместна тогда и только тогда, когда ранг основной матрицы равен рангу расширенной матрицы:
\[
\operatorname{rang} A = \operatorname{rang} \tilde{A}.
\]
\end{theorem}
